\documentclass[12pt,a4paper]{article}
\usepackage[utf8]{inputenc}
\usepackage[T2A]{fontenc}
\usepackage[russian]{babel}
\usepackage{amsmath, amssymb, amsthm}
\usepackage{geometry}
\geometry{top=2cm, bottom=2cm, left=2.5cm, right=1.5cm}
\usepackage{tikz}
\usetikzlibrary{automata, positioning, arrows.meta}
\usepackage{hyperref}

\theoremstyle{definition}
\newtheorem{definition}{Определение}
\theoremstyle{plain}
\newtheorem{theorem}{Теорема}
\newtheorem{lemma}{Лемма}
\theoremstyle{remark}
\newtheorem{example}{Пример}

\title{Билет №11 \\ \small Конечные автоматы и регулярные языки, их эквивалентность. Детерминизация и минимизация автоматов. (ИТМО) \\ Автоматы. Алгебры регулярных выражений. Теорема Клини о регулярных языках. (СПбГУ)}
\author{Сериков Владислав Александрович}
\date{21 апреля 2026}

\begin{document}

\maketitle
\tableofcontents
\newpage

\section{Языки, грамматики и распознаватели}

\begin{definition}[Грамматика и Язык]
\textbf{Порождающая грамматика} — это четверка $G = (V_T, V_N, P, S)$, где $V_T$ — алфавит терминальных символов, $V_N$ — алфавит нетерминальных символов, $P$ — конечное множество правил вывода (подстановок), $S \in V_N$ — начальный символ.
\textbf{Язык}, порождаемый грамматикой $G$, — это множество всех цепочек из терминальных символов, которые могут быть выведены из $S$: $L(G) = \{\alpha \in V_T^* \mid S \Rightarrow^* \alpha\}$.
\end{definition}

В иерархии Хомского грамматики типа 3 называются \textbf{регулярными} (праволинейными или леволинейными). Альтернативой порождающим грамматикам служат \textbf{распознаватели} — абстрактные машины, получающие на вход цепочку символов и принимающие решение: допустить её (цепочка принадлежит языку) или отвергнуть. Для регулярных языков таким распознавателем является конечный автомат.

\section{Основы теории конечных автоматов}

\begin{definition}[Автомат и Такт]
\textbf{Автомат} — это абстрактный дискретный исполнитель, преобразующий последовательность входных символов в выходные. Работа автомата происходит не непрерывно, а мгновенно в дискретные моменты времени, называемые \textbf{тактами}.
\end{definition}

\begin{definition}[Состояния автомата]
Значение выходного символа автомата зависит от входного символа и предыстории (предыдущей последовательности входов). Каждому подмножеству эквивалентных предысторий соответствует уникальное \textbf{состояние автомата}. Для хранения текущего состояния используется внутренняя память. Набор всех возможных состояний называется \textbf{множеством состояний $S$}.
Если множество $S$ конечно, автомат называется \textbf{конечным автоматом (КА)}.
\end{definition}

\begin{definition}[Детерминированный автомат-преобразователь]
Формально КА задается как шестерка $A = \langle X, Y, S, s_0, \delta, \lambda \rangle$, где $X$ и $Y$ — входной и выходной алфавиты, $s_0 \in S$ — начальное состояние, $\delta: S \times X \to S$ — функция переходов, $\lambda: S \times X \to Y$ — функция выходов.
Автомат \textbf{детерминированный}, если для любого состояния $s$ и входного символа $x$ переход $\delta(s, x)$ является единственным.
Такой автомат называется \textbf{преобразователем}, так как он строит выходную цепочку по входной.
\end{definition}

Способы задания автомата:
\begin{itemize}
    \item \textbf{Диаграмма переходов} — это направленный помеченный граф, вершины которого соответствуют состояниям, а дуги — переходам (помечаются входными и/или выходными символами).
    \item \textbf{Таблица функций переходов} — матрица (или таблица), строки которой соответствуют состояниям, а столбцы — входным символам. На пересечении указывается новое состояние.
\end{itemize}

\section{Автоматы-распознаватели и эквивалентность}

\begin{definition}[Конечный автомат-распознаватель]
Это ДКА «без выходов», определяемый пятеркой $A = \langle X, S, s_0, \delta, F \rangle$, где $F \subseteq S$ — множество допускающих (заключительных) состояний. Цепочка допускается, если после её прочтения автомат оказывается в состоянии $s \in F$.
\end{definition}

\begin{definition}[Эквивалентность автоматов]
Два автомата называются \textbf{эквивалентными}, если они допускают один и тот же язык.
\end{definition}

\begin{definition}[Недетерминированный конечный автомат и $\varepsilon$-переходы]
\textbf{Недетерминированный конечный автомат (НКА)} имеет функцию переходов вида $\delta: S \times X \to 2^S$ (возвращает множество возможных состояний).
Если функция переходов допускает вид $\delta: S \times (X \cup \{\varepsilon\}) \to 2^S$, говорят об НКА с \textbf{$\varepsilon$-переходами}. $\varepsilon$-переход — это смена состояния автомата без чтения входного символа. По теореме Рабина-Скотта классы языков ДКА и НКА совпадают.
\end{definition}

\section{Достижимость, эквивалентность состояний и минимизация}

\begin{definition}[Расширенная функция переходов и Достижимость]
\textbf{Расширенная функция переходов} $\delta^*: S \times X^* \to S$ определена на множестве входных цепочек (а не одиночных символов): $\delta^*(s, \varepsilon) = s$, $\delta^*(s, a\alpha) = \delta^*(\delta(s, a), \alpha)$. \\
Состояние $s$ называется \textbf{достижимым}, если существует цепочка $\alpha \in X^*$ такая, что $\delta^*(s_0, \alpha) = s$.
\end{definition}

\begin{definition}[Эквивалентность состояний и $k$-эквивалентность]
Два состояния $s_1$ и $s_2$ называются \textbf{эквивалентными} ($s_1 \equiv s_2$), если для любой цепочки $\alpha \in X^*$ выполняется: $\delta^*(s_1, \alpha) \in F \iff \delta^*(s_2, \alpha) \in F$. \\
Два состояния называются \textbf{$k$-эквивалентными}, если это условие выполняется для всех цепочек длины $|\alpha| \le k$.
\end{definition}

\begin{definition}[Минимальный ДКА]
ДКА без недостижимых и эквивалентных состояний является \textbf{минимальным ДКА}.
\end{definition}

\textit{Алгоритм минимизации (см. Теорему об эквивалентности):}
Основан на последовательном построении разбиений множества состояний $P_0, P_1, \dots, P_k$ на классы $k$-эквивалентности. $P_0$ делит состояния на $F$ и $S \setminus F$. На шаге $k$ состояния остаются в одном классе, если их переходы по всем символам $x \in X$ ведут в состояния, которые были $(k-1)$-эквивалентны (лежали в одном классе в $P_{k-1}$). Сходимость ($P_k = P_{k-1}$) дает состояния минимального ДКА.

\section{Регулярные множества, языки и выражения}

\begin{definition}[Регулярное множество и язык]
\textbf{Регулярное множество} над алфавитом $X$ определяется индуктивно:
\begin{enumerate}
    \item $\emptyset$, $\{\varepsilon\}$, $\{a\}$ (где $a \in X$) — регулярные множества.
    \item Если $P$ и $Q$ — регулярные множества, то их объединение ($P \cup Q$), конкатенация ($PQ$) и итерация/звезда Клини ($P^*$) также являются регулярными множествами.
\end{enumerate}
Язык, представляющий собой регулярное множество, называется \textbf{регулярным языком}.
\end{definition}

\begin{definition}[Регулярное выражение]
\textbf{Регулярное выражение} — это способ определения регулярного языка, альтернативный порождающим грамматикам. Оно строится из символов $\emptyset, \varepsilon, a$ с применением операций объединения ($p|q$), конкатенации ($pq$) и замыкания ($p^*$). В алгебре Клини эти выражения подчиняются аксиомам дистрибутивности, ассоциативности и идемпотентности.
\end{definition}

\textbf{Теорема Клини} утверждает эквивалентность всех трёх подходов: языки ДКА, языки регулярных грамматик и языки регулярных выражений — это в точности класс регулярных языков.

\section{Свойства регулярных языков и разрешимость}

\begin{definition}[Замкнутость]
Множество называется \textbf{замкнутым} относительно некоторой операции, если в результате её выполнения над элементами этого множества получается элемент, принадлежащий тому же множеству.
\end{definition}

Класс регулярных языков образует булеву алгебру и \textbf{замкнут относительно операций}:
\begin{itemize}
    \item Объединения, пересечения и разности;
    \item Дополнения (относительно $X^*$);
    \item Обращения (реверса цепочек);
    \item Итерации (звезды Клини) и конкатенации;
    \item Гомоморфизма (замены символов алфавита на цепочки).
\end{itemize}

\begin{definition}[Разрешимость проблем]
Для регулярных языков существует алгоритмическое решение (разрешимы) следующих фундаментальных проблем:
\begin{enumerate}
    \item \textbf{Проблема принадлежности цепочки языку:} Решается тривиальным запуском ДКА на данной цепочке.
    \item \textbf{Проблема пустоты языка ($L(G) = \emptyset$):} Сводится к проверке достижимости хотя бы одного заключительного состояния $s \in F$ из начального состояния $s_0$ в графе переходов автомата.
    \item \textbf{Проблема эквивалентности двух языков:} Разрешима путем построения минимальных ДКА для обоих языков. Если автоматы изоморфны, языки эквивалентны. (Либо через проверку пустоты симметрической разности языков: $(L_1 \cap \overline{L_2}) \cup (\overline{L_1} \cap L_2) = \emptyset$).
\end{enumerate}
\end{definition}

\section{Лемма о разрастании (накачке) для регулярных языков}

Лемма о разрастании (Pumping Lemma) является мощным инструментом для доказательства того, что некоторый язык \textbf{не является регулярным}. Она описывает необходимое (но не достаточное) условие регулярности языка.

\begin{lemma}[О разрастании / накачке]
Пусть $L$ — регулярный язык. Тогда существует константа $n > 0$ (зависящая только от языка $L$), такая, что любую цепочку $\alpha \in L$, длина которой $|\alpha| \ge n$, можно представить в виде конкатенации трёх подцепочек $\alpha = \beta\gamma\delta$, удовлетворяющих следующим условиям:
\begin{enumerate}
    \item $|\gamma| > 0$ (подцепочка $\gamma$ не пуста, $\gamma \neq \varepsilon$);
    \item $|\beta\gamma| \le n$ (длина префикса ограничена константой $n$);
    \item Для любого целого $i \ge 0$ цепочка $\beta\gamma^i\delta \in L$ (цепочка «разрастается» при повторении $\gamma$).
\end{enumerate}
\end{lemma}

\begin{proof}
Так как язык $L$ регулярен, существует минимальный детерминированный конечный автомат $A = \langle X, S, s_0, \delta, F \rangle$, допускающий $L$. Пусть $n = |S|$ — количество состояний этого автомата.

Рассмотрим цепочку $\alpha \in L$ такую, что $|\alpha| \ge n$. Обозначим её символы как $\alpha = x_1 x_2 \dots x_m$, где $m \ge n$. При чтении этой цепочки автомат проходит через последовательность состояний:
$p_0, p_1, p_2, \dots, p_m$,
где $p_0 = s_0$, и $p_k = \delta^*(s_0, x_1\dots x_k)$.

Рассмотрим первые $n+1$ состояний в этой последовательности: $p_0, p_1, \dots, p_n$. Поскольку всего состояний в автомате $n$, по \textbf{принципу Дирихле} (Pigeonhole principle), в этой последовательности хотя бы одно состояние должно повториться. Пусть это состояние $p_j = p_k$, где $0 \le j < k \le n$.

Тогда мы можем разбить цепочку $\alpha$ на три части:
\begin{itemize}
    \item $\beta = x_1 \dots x_j$ (переход от $p_0$ к $p_j$);
    \item $\gamma = x_{j+1} \dots x_k$ (цикл от $p_j$ до $p_k$, возвращающий в то же состояние);
    \item $\delta = x_{k+1} \dots x_m$ (переход от $p_k$ к конечному допускающему состоянию $p_m \in F$).
\end{itemize}
Проверим условия:
1. $k > j$, следовательно, длина $|\gamma| = k - j > 0$.
2. $k \le n$, следовательно, длина префикса $|\beta\gamma| = k \le n$.
3. Так как чтение подцепочки $\gamma$ возвращает автомат в то же самое состояние, автомат может пройти этот цикл $i$ раз (включая $0$ раз, т.е. пропустив его). При этом он всегда окажется в состоянии $p_j$, откуда по цепочке $\delta$ гарантированно дойдет до допускающего состояния. Следовательно, $\beta\gamma^i\delta \in L$ для любого $i \ge 0$. Лемма доказана.
\end{proof}

\begin{example}[Применение леммы о накачке]
Докажем, что язык $L = \{a^k b^k \mid k > 0\}$ \textbf{не является регулярным}.

Доказательство от противного. Допустим, язык $L$ регулярен. Тогда для него выполняется лемма о накачке и существует константа $n > 0$.

Выберем цепочку $\alpha = a^n b^n \in L$. Её длина $|\alpha| = 2n \ge n$. По лемме о накачке, её можно разбить на $\alpha = \beta\gamma\delta$, где $|\beta\gamma| \le n$ и $|\gamma| > 0$.

Так как префикс длины $n$ цепочки $\alpha$ состоит целиком из символов $a$, подцепочка $\beta\gamma$ также состоит только из символов $a$. Значит, $\gamma = a^m$, где $m > 0$.

По лемме, цепочка $\beta\gamma^0\delta$ (то есть $\beta\delta$) должна принадлежать $L$.
Вычислим её: при $i=0$ мы «вырезали» $m$ символов $a$. Новая цепочка имеет вид $a^{n-m} b^n$.
Поскольку $m > 0$, количество символов $a$ ($n-m$) не равно количеству символов $b$ ($n$). Следовательно, полученная цепочка \textbf{не принадлежит} языку $L$. Получено противоречие с леммой о накачке. Язык $L$ не является регулярным.
\end{example}

\newpage

\section{Теорема Клини: эквивалентность формализмов}

Теорема Клини объединяет грамматический (РВ), автоматный (КА) и алгебраический подходы к регулярным языкам.

\begin{theorem}[Теорема Клини]
Класс языков, распознаваемых конечными автоматами, в точности совпадает с классом языков, которые могут быть описаны регулярными выражениями.
\end{theorem}

Доказательство теоремы конструктивно и сводится к алгоритмам взаимных преобразований: \textbf{РВ $\to$ НКА $\to$ ДКА $\to$ РВ}. Переход НКА $\to$ ДКА был рассмотрен ранее (Алгоритм построения подмножеств). Рассмотрим оставшиеся два перехода.

\subsection{Построение НКА по регулярному выражению (Алгоритм Томпсона)}
Алгоритм Мак-Нотона — Ямады — Томпсона строит НКА с $\varepsilon$-переходами по структуре регулярного выражения (индукцией по длине РВ). Полученный НКА всегда имеет ровно одно начальное и ровно одно конечное состояние.

\textbf{Базис индукции:}
\begin{itemize}
    \item Для $\emptyset$: автомат состоит из состояний $s_0, s_f$ без переходов.
    \item Для $\varepsilon$: автомат состоит из $s_0, s_f$ с переходом $s_0 \xrightarrow{\varepsilon} s_f$.
    \item Для символа $a \in X$: автомат состоит из $s_0, s_f$ с переходом $s_0 \xrightarrow{a} s_f$.
\end{itemize}

\textbf{Индукционный переход (операции над НКА $M(p)$ и $M(q)$):}
\begin{enumerate}
    \item \textbf{Объединение $p | q$:} Вводятся новые $s_0, s_f$. Добавляются $\varepsilon$-переходы из $s_0$ в начальные состояния $M(p)$ и $M(q)$, и $\varepsilon$-переходы из конечных состояний $M(p)$ и $M(q)$ в $s_f$.
    \item \textbf{Конкатенация $pq$:} Конечное состояние $M(p)$ объединяется с начальным состоянием $M(q)$. Начальным становится $s_0$ автомата $M(p)$, конечным — $s_f$ автомата $M(q)$.
    \item \textbf{Звезда Клини $p^*$:} Вводятся новые $s_0, s_f$. Добавляются $\varepsilon$-переходы:
    \begin{itemize}
        \item Из $s_0$ в начальное состояние $M(p)$ и напрямую в $s_f$ (пропуск);
        \item Из конечного состояния $M(p)$ в начальное состояние $M(p)$ (повторение) и в $s_f$ (выход).
    \end{itemize}
\end{enumerate}

\subsection{Построение РВ по конечному автомату (Метод алгебраических уравнений)}
Задан ДКА (или НКА) $A = \langle X, S, s_0, \delta, F \rangle$. Требуется найти регулярное выражение для $L(A)$.

\textbf{Шаг 1.} Для каждого состояния $s_i \in S$ введем переменную $X_i$, обозначающую регулярное множество (язык), допускаемое автоматом, если бы $s_i$ было начальным состоянием.
Искомый язык всего автомата — это $X_0$ (переменная, соответствующая $s_0$).

\textbf{Шаг 2.} Составим систему уравнений. Для каждого $s_i$:
$$ X_i = \sum_{(s_i, a) \to s_j} a X_j + \begin{cases}
\varepsilon, & \text{если } s_i \in F \\
\emptyset, & \text{если } s_i \notin F
\end{cases} $$
Здесь знак $+$ означает объединение (или символ $|$).

\textbf{Шаг 3. Решение системы (Лемма Ардена).}
\begin{lemma}[Лемма Ардена]
Если уравнение имеет вид $X = \alpha X + \beta$, и $\varepsilon \notin L(\alpha)$, то уравнение имеет единственное решение: $X = \alpha^* \beta$.
\end{lemma}

Используя метод подстановки и лемму Ардена, мы последовательно исключаем переменные из системы, пока не получим явное выражение для $X_0$.

\begin{example}
Рассмотрим ДКА: $s_0 \xrightarrow{a} s_1$, $s_1 \xrightarrow{b} s_0$, $s_1 \xrightarrow{a} s_1$. Состояние $s_1 \in F$.
Система уравнений:
\begin{align*}
    X_0 &= a X_1 \\
    X_1 &= a X_1 + b X_0 + \varepsilon
\end{align*}
Подставим $X_0$ во второе уравнение:
$$ X_1 = a X_1 + b (a X_1) + \varepsilon = (a + ba) X_1 + \varepsilon $$
Применяем Лемму Ардена ($\alpha = a+ba$, $\beta = \varepsilon$):
$$ X_1 = (a + ba)^* \varepsilon = (a + ba)^* $$
Теперь найдем искомое $X_0$:
$$ X_0 = a X_1 = a (a + ba)^* $$
Полученное регулярное выражение: $a(a|ba)^*$.
\end{example}

\end{document}
